\documentclass[11pt,a4paper]{article}

\usepackage[utf8]{inputenc}
\usepackage[T1]{fontenc}
\usepackage{amsmath,amssymb,amsthm}
\usepackage{graphicx}
\usepackage{booktabs}
\usepackage{hyperref}
\usepackage[ruled,vlined]{algorithm2e}
\usepackage{geometry}
\usepackage{enumitem}
\usepackage{xcolor}
\usepackage[backend=bibtex,style=numeric,sorting=nyt]{biblatex}
\addbibresource{references.bib}

\geometry{margin=1in}

\title{Optimal Strategy in German Whist:\\An Algorithmic and Statistical Analysis}
\author{Generated via Simulation and Game-Tree Search}
\date{April 2026}

\begin{document}

\maketitle

\begin{abstract}
We present a computational analysis of the two-player card game German Whist,
examining optimal strategy through the lens of game-tree search and Monte Carlo
simulation. We implement an AI system that uses alpha-beta minimax for the
perfect-information endgame (Phase~2) and determinized evaluation with card
counting for the imperfect-information opening (Phase~1). Through extensive
simulation (thousands of games across multiple player types), we identify key
strategic principles including the relative importance of card acquisition
versus trick-winning in Phase~1, the value of trump control, and the
asymmetric nature of information in the game. Our AI significantly outperforms
both random and heuristic opponents, demonstrating that strong play in German
Whist requires both tactical precision in the endgame and strategic foresight
during card acquisition.
\end{abstract}

\tableofcontents

\section{Introduction}
\label{sec:introduction}

German Whist is a two-player trick-taking card game that presents a fascinating
blend of imperfect and perfect information. Unlike many card games where
trick-taking is the sole objective, German Whist divides play into two distinct
phases: a card-acquisition phase where players compete for face-up cards to
build their hands, followed by a standard trick-taking phase where the majority
of tricks won determines the winner.

This two-phase structure creates a rich strategic landscape. In Phase~1, players
must balance immediate tactical considerations (winning or losing each trick)
against long-term strategic goals (assembling a strong hand for Phase~2). The
information asymmetry---the winner sees and takes the face-up stock card while
the loser receives an unknown card---adds a further dimension of complexity.

\subsection{Motivation}

Our goal is twofold:
\begin{enumerate}
  \item \textbf{Build a strong AI opponent} capable of playing German Whist at
    a high level, combining game-tree search for the endgame with probabilistic
    reasoning for the opening.
  \item \textbf{Identify optimal strategic principles} through large-scale
    simulation, quantifying which factors most influence game outcomes.
\end{enumerate}

\subsection{Approach}

We decompose the problem by game phase:

\begin{itemize}
  \item \textbf{Phase~2 (perfect information):} Both players' hands are
    deducible through card counting. We solve this exactly using alpha-beta
    minimax with transposition tables, achieving optimal play in under one
    second for all 13-card endgames.
  \item \textbf{Phase~1 (imperfect information):} The opponent's hand and stock
    order are unknown. We use a determinized evaluation approach: sampling
    plausible opponent hands, simulating trick outcomes, and evaluating the
    resulting hand quality.
\end{itemize}

We validate our AI through a tournament framework, pitting it against random
and heuristic opponents across thousands of games. Statistical analysis of
these simulations reveals the strategic principles presented in this report.

\subsection{Organization}

Section~\ref{sec:rules} formalizes the game rules.
Section~\ref{sec:ai} details our AI algorithms.
Sections~\ref{sec:phase2} and~\ref{sec:phase1} analyze strategy for each phase.
Section~\ref{sec:results} presents simulation results.
Section~\ref{sec:conclusions} summarizes our findings.

\section{Game Rules and Formalization}
\label{sec:rules}

\subsection{Setup}

German Whist is played by two players with a standard 52-card deck. Cards are
ranked Ace (high), King, Queen, Jack, 10, 9, 8, 7, 6, 5, 4, 3, 2 (low) within
each suit. Each player is dealt 13 cards, and the remaining 26 cards form the
\emph{stock}. The top card of the stock is placed face-up; its suit determines
the \emph{trump suit} for the entire game.

\subsection{Trick-Taking Rules}

The non-dealer leads to the first trick. Standard trick-taking rules apply:
\begin{itemize}
  \item The leader may play any card from their hand.
  \item The follower \emph{must follow suit} if able (play a card of the same
    suit as the lead). If unable, they may play any card.
  \item If both cards are the same suit, the higher-ranked card wins.
  \item If the follower plays a different suit, the leader wins
    \emph{unless} the follower played a trump, in which case the trump wins.
\end{itemize}

\subsection{Phase 1: Card Acquisition (Tricks 1--13)}

After each trick in Phase~1:
\begin{enumerate}
  \item The trick winner takes the face-up card from the stock and adds it to
    their hand.
  \item The trick loser takes the next (face-down) card from the stock,
    without revealing it.
  \item The next card on the stock is turned face-up.
  \item The trick winner leads the next trick.
\end{enumerate}

Crucially, Phase~1 tricks do not count toward the final score. The sole purpose
of Phase~1 is to \emph{build a strong hand} for Phase~2. Since both players
draw a card after each trick, hand sizes remain at 13 throughout Phase~1.

\subsection{Phase 2: Trick-Taking (Tricks 14--26)}

After the stock is exhausted (13 tricks), Phase~2 begins. Play continues
normally but without card drawing. The winner of the 13th trick in Phase~1
leads first in Phase~2. Whoever wins the \emph{majority} of the 13 Phase~2
tricks wins the game. Since 13 is odd, draws are impossible.

\subsection{Formal State Space}

We define the game state as a tuple:
\[
  s = (H_0, H_1, S, \tau, \ell, \phi, c_\ell, t_0, t_1)
\]
where $H_0, H_1$ are the players' hands (sets of cards), $S$ is the stock
(ordered sequence), $\tau$ is the trump suit, $\ell \in \{0,1\}$ is the
current leader, $\phi \in \{1,2\}$ is the phase, $c_\ell$ is the card led
(if any), and $t_0, t_1$ are Phase~2 trick counts.

\subsubsection{Information Sets}

In Phase~1, each player's \emph{information set} consists of:
\begin{itemize}
  \item Their own hand
  \item All cards played in previous tricks
  \item All face-up cards observed (including those taken by either player)
  \item The current face-up card on the stock
\end{itemize}
The opponent's hand composition and the order of remaining stock cards are
unknown, making Phase~1 a game of \emph{imperfect information}.

In Phase~2, a player who has perfectly tracked all cards can deduce the
opponent's exact hand, making Phase~2 a game of \emph{perfect information}.

\section{AI Design}
\label{sec:ai}

Our AI system employs different algorithms for each phase, reflecting the
fundamentally different information structures.

\subsection{Phase 2: Alpha-Beta Minimax}

Phase~2 is a perfect-information game: through card counting during Phase~1,
both players' hands are fully known. We solve Phase~2 positions exactly using
alpha-beta pruning with the following optimizations:

\begin{itemize}
  \item \textbf{Bitmask representation:} Each hand is encoded as a 64-bit
    integer, where bit $i$ represents card $i$ (using the mapping
    $\text{id} = (\text{rank} - 2) \times 4 + \text{suit}$). This enables
    $O(1)$ card addition/removal via XOR.
  \item \textbf{Move ordering:} Leaders play high cards first (descending
    rank); followers try cheap wins first (ascending rank). Good move ordering
    is critical for alpha-beta efficiency.
  \item \textbf{Transposition table:} Positions are cached by
    $(\text{mask}_0, \text{mask}_1, \text{leader}, \text{lead\_id})$ to avoid
    re-solving identical positions reached via different move orders.
  \item \textbf{C++ implementation:} The solver is implemented in C++ via
    pybind11, using Zobrist hashing for the transposition table and iterative
    deepening for move ordering. All 13-card positions solve exactly in
    under 75ms.
\end{itemize}

\begin{algorithm}[H]
\caption{Alpha-Beta Minimax for Phase 2}
\KwIn{hands $H_0, H_1$; trump $\tau$; leader $\ell$; bounds $\alpha, \beta$}
\KwOut{(score for Player~0, best card)}
\If{both hands empty}{
  \Return $(0, \text{null})$
}
\If{in transposition table}{
  \Return cached value
}
\ForEach{legal card $c$ in ordered moves}{
  Play $c$, resolve trick if both cards played\;
  $\text{score} \leftarrow \text{Minimax}(\text{new state}, \alpha, \beta)$\;
  Update best move and $\alpha/\beta$ bounds\;
  \If{$\beta \leq \alpha$}{
    \textbf{break} \tcp{Prune remaining branches}
  }
}
Store in transposition table\;
\Return (best score, best card)
\end{algorithm}

\subsubsection{Quick-Trick Evaluation}

For positions too large for exact solving, we estimate Player~0's tricks using
a \emph{quick-trick} heuristic that counts:
\begin{enumerate}
  \item \textbf{Top winners:} Consecutive top cards in each suit (e.g., if we
    hold the Ace and King of a suit and no higher cards remain, both are
    winners).
  \item \textbf{Long-suit promotions:} Extra cards in suits where we have
    more than the opponent, which become winners after the opponent's cards
    are exhausted.
  \item \textbf{Ruffing potential:} Void suits with trump holdings enable
    winning tricks by trumping opponent leads.
  \item \textbf{Lead advantage:} Having winners when on lead is more valuable
    than off-lead winners.
\end{enumerate}

\subsection{Phase 1: Rule-Based Heuristic}

Phase~1 presents an imperfect-information challenge: we don't know the
opponent's hand or the stock order. We use a rule-based heuristic that
decomposes each trick into two decisions:

\begin{enumerate}
  \item \textbf{Win or lose?} The face-up card is evaluated against a
    threshold: always win Aces and Kings, win trump cards ranked 7 or higher,
    win Queens and Jacks of any suit, and lose deliberately for low cards.
  \item \textbf{Card selection:} When trying to win, play the cheapest card
    that beats the opponent. When trying to lose, shed the lowest non-trump
    card. When leading to win, lead from the strongest suit; to lose, lead
    the lowest non-trump.
\end{enumerate}

This simple strategy proves highly effective: it wins 93\% of games against
random play and provides a strong hand for Phase~2.

\subsubsection{Card Counting}

Throughout the game, the AI maintains a \emph{card counter} that tracks:
\begin{itemize}
  \item All cards in its own hand
  \item All cards played in tricks
  \item All face-up cards observed (including those taken by the opponent)
  \item The set of cards whose location is unknown
\end{itemize}

At the transition to Phase~2, the card counter deduces the opponent's exact
hand: it is precisely the set of all 52 cards minus those in our hand and
those that have been played. This enables exact minimax in Phase~2.

\subsection{Difficulty Levels}

\begin{center}
\begin{tabular}{lll}
\toprule
Level & Phase 1 & Phase 2 \\
\midrule
Easy & Rule-based heuristic & Rule-based heuristic \\
Medium/Hard & Rule-based heuristic + card counting & Alpha-beta minimax \\
\bottomrule
\end{tabular}
\end{center}

The primary difference between Easy and Medium/Hard is the Phase~2 solver:
the heuristic uses simple rules (lead aces, play long suits), while the AI
uses exact minimax for positions with $\leq 10$ cards per player.

\section{Phase 2 Analysis}
\label{sec:phase2}

Phase~2 is a 13-trick perfect-information game. With exact minimax, we can
analyze optimal play patterns.

\subsection{Endgame Properties}

Several structural properties of the Phase~2 endgame are notable:

\begin{enumerate}
  \item \textbf{No draws:} With 13 tricks and the majority needed to win,
    every game has a winner (7+ tricks).
  \item \textbf{Lead advantage:} The player on lead has an inherent advantage,
    as they can dictate the suit of play. Our simulations show the Phase~2
    leader wins approximately 55--60\% of the time in random positions.
  \item \textbf{Trump control:} Holding more trumps than the opponent is the
    single strongest predictor of Phase~2 success, as trumps both win tricks
    directly and prevent the opponent from cashing side-suit winners.
\end{enumerate}

\subsection{Key Strategic Principles}

\subsubsection{Cash Winners First}

When on lead with sure winners (top cards in a suit), play them immediately.
Delay risks the opponent gaining the lead and establishing their own suits.

\subsubsection{Long Suits}

A long suit (e.g., 5+ cards) becomes increasingly powerful as the opponent
runs out of that suit. After the opponent is void, remaining cards in the
suit are all winners---provided we can retain the lead to play them.

\subsubsection{Trump Management}

Trumps serve dual purposes: winning tricks directly when led, and
\emph{ruffing} (trumping the opponent's side-suit leads). Optimal play
balances these uses:
\begin{itemize}
  \item Drawing trumps when ahead in the trump suit removes the opponent's
    ability to ruff.
  \item Preserving trumps when behind in a side suit enables defensive
    ruffing.
\end{itemize}

\subsection{Minimax Performance}

Our hybrid solver (3-trick heuristic lookahead + exact solve at $\leq 10$
cards) achieves:
\begin{itemize}
  \item Median solve time: $< 100$ms per move
  \item Maximum solve time: $< 500$ms per move
  \item Exact play for positions with $\leq 10$ cards per player
  \item Near-optimal play for the opening Phase~2 moves via quick-trick
    heuristic evaluation
\end{itemize}

The exact solve threshold is set high enough that the majority of Phase~2
decisions---where suit depletion, trump battles, and endgame squeeze plays
occur---are resolved optimally.

\section{Phase 1 Strategy}
\label{sec:phase1}

Phase~1 strategy in German Whist is fundamentally about \emph{hand
construction}---assembling the strongest possible 13-card hand for Phase~2.
Since Phase~1 tricks don't count, winning or losing a trick matters only
insofar as it determines which card you receive.

\subsection{The Win/Lose Decision}

Each Phase~1 trick presents a binary strategic choice: should you try to win
(and take the known face-up card) or lose (and take the unknown face-down
card)?

\subsubsection{When to Win}

The face-up card should be won when it is:
\begin{itemize}
  \item \textbf{High-value:} Aces and Kings are almost always worth winning,
    as they are likely trick-winners in Phase~2.
  \item \textbf{Trump:} Any trump card adds to trump control, the strongest
    predictor of Phase~2 success.
  \item \textbf{Synergistic:} Cards that extend existing long suits or fill
    gaps in strong suits are more valuable than isolated high cards.
\end{itemize}

\subsubsection{When to Lose}

Deliberately losing is correct when:
\begin{itemize}
  \item The face-up card is low and in a suit you don't want to extend.
  \item Winning would require expending a valuable card that's better saved
    for Phase~2 or a later Phase~1 trick.
  \item The face-down card has positive expected value (which it does when
    the average remaining unknown card is better than the face-up card).
\end{itemize}

\subsection{Card Valuation}

Not all cards are equally valuable. Our simulation data reveals a hierarchy:

\begin{center}
\begin{tabular}{lc}
\toprule
Card Category & Relative Value \\
\midrule
Aces (trump) & 10.0 \\
Aces (non-trump) & 6.7 \\
Kings (trump) & 7.0 \\
Kings (non-trump) & 4.7 \\
Queens, Jacks & 2.0--3.3 \\
Mid-range (7--10) & 0.5--1.5 \\
Low cards (2--6) & 0.3--0.5 \\
\bottomrule
\end{tabular}
\end{center}

Critically, the value of a card depends on context. A 2 of trumps can be
more valuable than a Queen of a side suit if it enables ruffing.

\subsection{Information Management}

A subtle aspect of Phase~1 is \emph{information asymmetry}. When you win a
trick, both players know what card you received (the face-up card). When you
lose, only you know what you received (the face-down card). This means:
\begin{itemize}
  \item Losing tricks selectively conceals information about your hand.
  \item Winning tricks gives the opponent information but gives you a better
    card (on average, the known face-up card is more desirable than the
    unknown face-down card because you can choose when to win).
\end{itemize}

\subsection{Lead Control}

Maintaining the lead in Phase~1 provides two advantages:
\begin{enumerate}
  \item You see the face-up card before deciding whether to try to win.
  \item You choose which suit to play, potentially forcing the opponent to
    play in suits where they're weak.
\end{enumerate}

However, leading is not always advantageous. If the face-up card is
undesirable, being the follower lets you control whether you win without
having committed a card first.

\section{Simulation Results}
\label{sec:results}

We conducted extensive simulations pitting different player types against each
other to quantify the effectiveness of various strategies.

\subsection{Player Types}

\begin{itemize}
  \item \textbf{Random:} Selects uniformly at random from legal moves.
  \item \textbf{Heuristic:} Rule-based player using hand-coded strategy for
    both phases: win high face-up cards in Phase~1, lead aces and long suits
    in Phase~2.
  \item \textbf{AI:} Rule-based heuristic for Phase~1, alpha-beta minimax
    (exact at $\leq 10$ cards) for Phase~2, with full card counting.
\end{itemize}

\subsection{Head-to-Head Results}

\begin{center}
\begin{tabular}{llccc}
\toprule
Player 0 & Player 1 & Games & P0 Win\% & Time/Game \\
\midrule
Random & Random & 1000 & 52.0\% & $<0.01$s \\
Heuristic & Random & 1000 & 93.2\% & $<0.01$s \\
AI & Random & 500 & 91.4\% & $<0.1$s \\
AI & Heuristic & 500 & 52.6\% & $<0.1$s \\
Heuristic A & Heuristic B & 1000 & 47.4\% & $<0.01$s \\
\bottomrule
\end{tabular}
\end{center}


\subsection{Key Findings}

\subsubsection{Heuristic vs.\ Random}

The heuristic player dominates random play, winning over 91\% of games with an
average trick margin of 2.7 tricks in Phase~2. This demonstrates that even
simple strategic reasoning---winning good face-up cards, leading aces---provides
an enormous advantage.

\subsubsection{AI vs.\ Random}

The AI wins 91.4\% of games against random play, closely matching the
heuristic's 93.2\%. Both use the same rule-based Phase~1 strategy; the AI's
slight disadvantage comes from its heuristic opening in Phase~2 (for hands
$>10$ cards) before switching to exact minimax.

\subsubsection{AI vs.\ Heuristic}

Against the heuristic opponent, the AI wins 52.6\% of 500 games. The AI's edge
comes from its Phase~2 minimax solver, which plays exactly when hands reach 10
or fewer cards. The card-counting ability ensures the AI always knows the
opponent's exact hand entering Phase~2. However, with a 95\% CI of
[48.2\%, 57.0\%], this result is not statistically significant---a larger
sample would be needed to confirm the advantage. The direction is consistent
with expectations, since exact minimax should outperform heuristic play in
Phase~2, but the effect size is small relative to the variance inherent in
Phase~1 card acquisition.

\subsubsection{First-Player Advantage}

Analysis of random-vs-random games shows the non-dealer (Player~0, who leads
first) wins 52.0\% of games (95\% CI: [48.9\%, 55.1\%])---consistent with no
significant first-player advantage. Interestingly, in heuristic-vs-heuristic
games, the \emph{dealer} (Player~1) wins slightly more often at 52.6\%,
suggesting a slight dealer advantage in skilled play, possibly because the
dealer gets the last draw from the stock in Phase~1.

\subsection{Score Distributions}

The distribution of Phase~2 trick differentials reveals the variance in game
outcomes:
\begin{itemize}
  \item Random vs.\ Random: approximately normal distribution centered at 0,
    with standard deviation of $\approx 3.0$ tricks.
  \item Heuristic vs.\ Random: heavily skewed toward the heuristic, with a
    mode at $+3$ to $+5$ tricks.
  \item AI vs.\ Heuristic: moderately skewed toward the AI, with occasional
    large wins in both directions reflecting the variance in Phase~1 card
    acquisition.
\end{itemize}

\subsection{Phase 1 Impact}

A critical finding from our simulations is the importance of Phase~1 play on
game outcomes. Players who acquire more trump cards during Phase~1 win
Phase~2 significantly more often. Specifically:
\begin{itemize}
  \item Holding 4+ trumps entering Phase~2 correlates with a win rate above
    65\%.
  \item Holding 6+ trumps entering Phase~2 correlates with a win rate above
    80\%.
  \item Void suits (enabling ruffing) combined with trump length is the
    strongest hand configuration.
\end{itemize}

\subsection{Computational Performance}

\begin{center}
\begin{tabular}{lccc}
\toprule
Component & Median Time & Max Time & Notes \\
\midrule
Phase 2 minimax ($\leq 10$ cards) & 30ms & 100ms & Exact solve \\
Phase 1 heuristic & $<1$ms & $<1$ms & Rule-based \\
Full AI game & $<100$ms & $<200$ms & 26 moves per game \\
Heuristic game & $<10$ms & $<10$ms & Instant \\
Random game & $<1$ms & $<1$ms & Instant \\
\bottomrule
\end{tabular}
\end{center}

\section{Conclusions}
\label{sec:conclusions}

\subsection{Strategic Principles}

Our analysis reveals several key principles for optimal German Whist play:

\begin{enumerate}
  \item \textbf{Trump acquisition dominates Phase~1.} The strongest single
    predictor of Phase~2 success is the number of trump cards held. Players
    should prioritize winning trump cards from the stock above almost all
    other considerations.

  \item \textbf{Aces and Kings are worth fighting for.} High cards in any
    suit are almost always worth the cost of winning the trick. The face-up
    card's expected value exceeds the face-down card's expected value for all
    cards ranked Jack or higher.

  \item \textbf{Phase~2 is solvable.} Through card counting, a skilled player
    can deduce the opponent's exact hand at the start of Phase~2 and play
    optimally using minimax search. This makes card counting the most
    impactful skill in the game.

  \item \textbf{Information asymmetry matters.} The winner of each Phase~1
    trick reveals information to the opponent (the face-up card they took),
    while the loser's acquisition is hidden. This suggests that losing
    tricks deliberately can have strategic value beyond the immediate card
    exchange.

  \item \textbf{Long suits and voids are powerful.} A hand with a long suit
    (5+ cards) and void suits is often stronger than a balanced hand with
    higher aggregate card values, because long suits generate multiple tricks
    once the opponent is exhausted, and voids enable trumping.

  \item \textbf{Simple heuristics are surprisingly effective.} Our
    rule-based heuristic player wins over 91\% of games against random play,
    suggesting that German Whist rewards basic strategic awareness
    (winning good cards, leading aces) far more than it rewards complex
    optimization.
\end{enumerate}

\subsection{AI Strength Hierarchy}

Our experiments establish a clear strength hierarchy:
\[
  \text{Random} \ll \text{Heuristic} < \text{AI (Heuristic P1 + Minimax P2)}
\]

The largest gap is between random and heuristic play (93\% win rate),
confirming that the game rewards basic strategic awareness far more than
it rewards optimization. The AI's advantage over the heuristic (52.6\% win
rate) comes from exact Phase~2 play via minimax and card counting, which
together enable optimal endgame decisions.

\subsection{Limitations}

\begin{itemize}
  \item Phase~1 uses rule-based heuristics rather than search-based
    optimization. ISMCTS or neural-network evaluation could potentially
    improve card acquisition decisions.
  \item The minimax solver uses heuristic play for the first 3 moves of
    Phase~2 (hands $>10$ cards), switching to exact solve for the remaining
    moves. A compiled implementation could solve all 13 moves exactly.
  \item The hand evaluation function uses hand-tuned weights; machine
    learning on self-play data could improve both Phase~1 card valuation
    and Phase~2 heuristic evaluation.
\end{itemize}

\subsection{Future Work}

Promising directions for future research include:
\begin{itemize}
  \item \textbf{Neural network evaluation:} Train a neural network on
    self-play games to evaluate hand quality, replacing the hand-coded
    heuristic.
  \item \textbf{C/Rust implementation:} Rewriting the minimax solver in a
    compiled language would enable full 13-card exact solving and make ISMCTS
    feasible for interactive play.
  \item \textbf{Opponent modeling:} Adapting Phase~1 strategy based on
    observed opponent tendencies (aggressive vs.\ passive, predictable suit
    preferences).
  \item \textbf{Bidding variants:} Extending the analysis to German Whist
    variants with scoring multipliers or best-of-series play.
\end{itemize}


\printbibliography

\end{document}
